\section{Strongly Correlated Regime: Frustration-Stabilized Composites}
  \label{sec:strongly-correlated-regime}

  \subsection{Real-Space Frustration and Raccordement Cost}
    \label{subsec:frustration-raccordement}

    In strongly correlated materials such as cuprates,
    the normal state cannot be described as a weakly interacting Fermi liquid.
    At low doping, the system is proximate to a Mott insulating phase
    with robust antiferromagnetic correlations.
    Within the present framework,
    this regime is characterized by a high density of
    real-space admissibility constraints on the fiber.

    We denote by $\mathcal{F}$ the local frustration density,
    defined operationally as a proxy for the fraction of
    nearest-neighbor raccordement constraints
    that cannot be simultaneously satisfied.
    A quantitative microscopic definition and its experimental grounding
    are provided in Appendix~\ref{app:frustration}.
    In cuprates, this quantity is tied to the
    antiferromagnetic structure factor $S(\pi,\pi)$,
    measurable by neutron scattering or RIXS,
    and we take
    \begin{equation}
      \mathcal{F}
      \propto
      \frac{S(\pi,\pi)}{S_{\max}},
      \label{eq:frustration-proxy}
    \end{equation}
    where $S_{\max}$ denotes the fully ordered reference value.

    In this regime,
    isolated $w=1$ excitations carry a high raccordement cost,
    as their presence locally increases
    the incompatibility of neighboring fiber phases.
    A bound $w=2$ composite,
    however,
    can reduce the global frustration
    by redistributing the phase mismatch
    over a larger region.
    Pair formation is therefore not driven by an external mediator,
    but by the minimization of the real-space fiber incompatibility.

    The projective gap $\Delta_\Pi^{\mathrm{amplitude}}$
    in this regime is set by the dominant frustration scale,
    which in cuprates is the magnetic exchange energy $J$.
    To leading order,
    \begin{equation}
      \Delta_\Pi^{\mathrm{amplitude}}
      \sim
      J .
      \label{eq:gap-frustration}
    \end{equation}

  \subsection{Symmetry Selection from Lattice Constraints}
    \label{subsec:symmetry-selection}

    The symmetry of the superconducting gap
    is determined by minimizing the total raccordement cost
    of the composite configuration
    under lattice symmetry constraints.

    In cuprates,
    the underlying lattice symmetry is $D_{4h}$,
    and the dominant frustration vector
    is $(\pi,\pi)$.
    The staggered nature of this frustration
    imposes a sign-alternating constraint
    on the fiber phase between neighboring plaquettes.

    The admissible composite must satisfy two conditions:
    \begin{enumerate}
      \item minimize overlap with the frustrated $(\pi,\pi)$ sector,
      \item preserve global fiber continuity without introducing defects.
    \end{enumerate}

    Among the irreducible representations of $D_{4h}$
    available to a $w=2$ composite,
    the lowest-cost solution under these constraints
    lies in the $B_{1g}$ sector.
    Its nodal structure aligns with the
    $(\pi,0)$ and $(0,\pi)$ directions,
    distributing the required sign alternation
    across the Fermi surface
    while avoiding localized raccordement singularities.
    The explicit mapping between the projective spectrum
    $\Delta_\Pi(\mathbf{k})$ and the ARPES gap structure
    underlying this nodal behavior
    is derived in Appendix~\ref{app:mapping}.

    Extended $s$-wave configurations
    do not simultaneously satisfy both conditions.
    Although certain $A_{1g}$ form factors change sign,
    they either accumulate excessive cost
    near the $(\pi,\pi)$ sector
    or require additional fiber defects
    to maintain global continuity.

    The selection of $d_{x^2-y^2}$ symmetry
    thus follows from a real-space geometric minimization problem,
    rather than from a momentum-space interaction kernel.

  \subsection{Amplitude--Phase Separation and the Pseudogap}
    \label{subsec:amplitude-phase}

    A crucial distinction between conventional and strongly correlated regimes
    is the separation of amplitude and phase scales.

    The formation of local $w=2$ composites
    occurs when
    \begin{equation}
      k_B T
      \lesssim
      \Delta_\Pi^{\mathrm{amplitude}}
      \sim
      J .
    \end{equation}
    We identify this scale with the pseudogap temperature $T^*$.

    However,
    global superconductivity requires
    phase locking across the system.
    The relevant quantity is the phase stiffness $\rho_s$,
    which measures the energetic cost
    of spatial variation of the fiber phase.

    In the doped regime,
    the density of active composites
    is proportional to the carrier concentration $\delta$.
    The stiffness therefore scales as
    \begin{equation}
      \rho_s^{(0)}
      \propto
      \delta \,
      \Delta_\Pi^{\mathrm{amplitude}}
      \sim
      \delta J .
      \label{eq:stiffness-scaling}
    \end{equation}

    The superconducting transition temperature
    is governed by phase coherence rather than amplitude formation.
    In quasi-two-dimensional systems,
    the transition is controlled by vortex unbinding,
    leading to the Berezinskii--Kosterlitz--Thouless criterion
    \begin{equation}
      k_B T_c
      \simeq
      \frac{\pi}{2}
      \rho_s(T_c^-)
      \sim
      \frac{\pi}{2}
      \delta J ,
      \label{eq:tc-deltaJ}
    \end{equation}
    up to corrections from interlayer coupling.

    This naturally explains the separation
    between $T^*$ and $T_c$,
    as well as the linear scaling of $\rho_s$
    with doping observed experimentally
    (Uemura relation).

  \subsection{Dome Structure and Frustration Ratio}
    \label{subsec:dome-structure}

    The superconducting dome arises
    from the competition between
    frustration amplitude and phase stiffness.

    At very low doping,
    $\mathcal{F}$ is large,
    but $\delta$ is small,
    yielding strong amplitude formation
    yet insufficient phase stiffness.
    The system remains antiferromagnetic.

    At optimal doping,
    $\mathcal{F}$ is reduced
    while $\delta$ is sufficient
    to sustain large $\rho_s$,
    maximizing $T_c$.

    In the overdoped regime,
    $\mathcal{F}$ becomes small,
    reducing the intrinsic stabilization
    of the composite,
    and $T_c$ decreases accordingly.

    The ratio
    \begin{equation}
      r_{\mathcal{F}}
      \equiv
      \text{first-order proxy for relative frustration amplitude}
    \end{equation}
    will play a central role
    in comparing material families,
    as developed in Section~\ref{sec:nickelates}.
